\documentclass[11pt]{article}
\usepackage[review]{acl}

\usepackage{times}
\usepackage{latexsym}
\usepackage[T1]{fontenc}
\usepackage[utf8]{inputenc}
\usepackage{microtype}
\usepackage{inconsolata}
\usepackage{graphicx}
\usepackage{booktabs}
\usepackage{multirow}
\usepackage{amsmath}
\usepackage{amssymb}
\usepackage{algorithm}
\usepackage{algpseudocode}
\usepackage{enumitem}

\title{StepKV: Step-Aware KV Cache Compression for LLM Agent Long-Horizon Reasoning}

\author{
First Author \\
Affiliation / Address line 1 \\
Affiliation / Address line 2 \\
\texttt{email@domain}
\And
Second Author \\
Affiliation / Address line 1 \\
Affiliation / Address line 2 \\
\texttt{email@domain}
}

\begin{document}
\maketitle

\begin{abstract}
Large language model (LLM) agents based on ReAct repeatedly generate \textit{Thought--Action--Observation} trajectories, causing the key--value (KV) cache to grow linearly with the number of reasoning steps.
Existing KV pruning methods such as H2O and SnapKV retain tokens by local attention saliency and therefore assume temporal locality.
We find that this assumption breaks in multi-step agentic reasoning: information introduced in an early step often becomes critical only several steps later (\emph{delayed reuse}), so token-only pruning evicts globally important context and causes severe accuracy collapse under limited KV budgets (e.g., H2O reaches 12.80 EM on HotpotQA at 50\% budget vs.\ 28.40 EM for full KV).

We propose \textbf{StepKV}, a step-aware KV cache compression framework that groups trajectory tokens by reasoning step and assigns each step a lightweight utility score from external trajectory signals---progress quality, novelty, repetition, and delayed citation.
At every step boundary, StepKV executes a three-phase scoring pipeline: (i) initialize the score of the newly finalized step; (ii) globally update older step scores when later text reuses their anchors; and (iii) combine token-level heavy-hitter scores with step utility and retain the top-$B$ tokens under budget $B=\rho n$.
Experiments on HotpotQA, 2WikiMultihopQA, MuSiQue, and BrowseComp with Qwen2.5-7B show that StepKV substantially outperforms H2O and SnapKV at matched budgets (26.50/35.50 EM/F1 vs.\ 12.80/20.24 for H2O on HotpotQA at 50\% budget) while closely approaching full-cache performance.
\end{abstract}

\section{Introduction}

Large language models (LLMs) have shown strong performance on complex reasoning when coupled with agent-style frameworks such as ReAct~\cite{yao2023react}, which interleave internal reasoning (\textit{Thought}), tool calls (\textit{Action}), and environment feedback (\textit{Observation}).
This paradigm enables multi-hop question answering over external corpora and has become a standard recipe for retrieval-augmented agents.

A practical bottleneck is the growing KV cache.
In ReAct, each step appends hundreds of new tokens to the decode cache; over a 7-step trajectory the cache can reach several thousand tokens.
Maintaining the full cache preserves accuracy but increases GPU memory and decode latency.
Recent token-level pruning methods such as H2O~\cite{zhang2023h2o} and SnapKV~\cite{li2024snapkv} reduce cache size by keeping attention-heavy tokens, and work well in standard autoregressive generation.

However, we observe that these methods fail dramatically in ReAct-style workloads.
On HotpotQA with BM25 retrieval over full Wikipedia, H2O at a 50\% KV budget drops to 12.80 EM, while a ReAct baseline without aggressive pruning reaches 28.00 EM.
The failure mode is systematic rather than random: early observations that appear locally unimportant are evicted, yet their entities reappear in later thoughts or actions.

We attribute this to \textbf{delayed reuse}.
Token importance in agent trajectories is not purely temporally local; it is \emph{step-conditional}.
An observation token may receive low attention immediately after generation, but the step that produced it can become semantically critical once later reasoning cites the same entities.
Token-only scores cannot represent this cross-step dependency.

We propose \textbf{StepKV}, which introduces explicit step structure into KV retention.
StepKV partitions the prunable cache region into reasoning steps, estimates a step utility $U_s$ from lightweight external signals (success, novelty, repetition, citation), and forms a hybrid token score $C_i=\alpha T_i+\beta U_{\mathrm{step}(i)}$ that combines H2O token saliency $T_i$ with step utility.
Crucially, step utility is maintained through a \textbf{three-phase update} at each step boundary rather than a one-shot local heuristic.

Our contributions are:
\begin{itemize}
    \item We systematically study KV cache pruning under ReAct workloads and identify delayed reuse as a key failure mode of token-only methods.
    \item We propose StepKV, a step-aware compression framework with anchor-based step utility and a three-phase scoring pipeline at step boundaries.
    \item We introduce a hybrid retention rule $C=\alpha T+\beta U$ under global budget $B=\rho n$, preserving both local token saliency and trajectory-level importance.
    \item Experiments on four multi-step QA benchmarks show that StepKV near-matches full KV while outperforming H2O and SnapKV by large margins under equal budgets; analyses confirm the value of citation and step utility.
\end{itemize}

\section{Related Work}

\subsection{Architectural Optimizations}

Architectural innovations reduce the intrinsic growth rate of KV storage.
Grouped-Query Attention (GQA)~\cite{gqa} and Multi-Head Latent Attention (MLA)~\cite{mla} share key/value heads across queries.
State-space and hybrid architectures~\cite{gu2023mamba,yang2024gateddelta,qwen2025next,nvidia2025nemotron} further alter memory scaling.
These methods lower per-token storage cost but do not decide \emph{which historical trajectory tokens to keep} once a long ReAct episode is underway.

\subsection{Sparse and Selective Attention}

Methods such as DeepSeek Sparse Attention (DSA)~\cite{liu2025dsa} reduce attention compute by selecting relevant keys at inference time.
They primarily address bandwidth and FLOPs rather than the physical size of the stored KV cache.
StepKV is complementary: it explicitly compresses stored trajectory context in agent settings.

\subsection{KV Cache Management and Compression}

StreamingLLM and H2O~\cite{zhang2023h2o} retain attention sinks or heavy-hitter tokens.
SnapKV~\cite{li2024snapkv} clusters local keys within attention windows; Scope~\cite{scope} separates prompt and generation pools.
Reasoning-oriented extensions include R-KV~\cite{cai2025rkv}, RaaS~\cite{hu2025raas}, LazyEviction~\cite{zhang2025lazyeviction}, ThinKV~\cite{ramachandran2025thinkv}, and SideQuest~\cite{sidequest}.
Most rely on token-level saliency and do not model step-wise structure or delayed cross-step reuse.
StepKV differs by grouping tokens by ReAct step and scoring steps with external trajectory signals before pruning.

\section{Methodology: Step-Aware KV Cache Compression}
\label{sec:methodology}

StepKV is invoked after each completed reasoning step to select a budgeted subset of the prunable KV cache.
At pruning point $t$, let $\mathcal{T}^{(t)}=\{1,\ldots,N_t\}$ denote prunable token indices in the trajectory region, and let $\mathcal{S}^{(t)}=\{1,\ldots,t{-}1\}$ denote finalized steps whose spans are registered.
Given keep ratio $\rho$, StepKV returns $\mathcal{K}^{(t)}\subseteq\mathcal{T}^{(t)}$ with $|\mathcal{K}^{(t)}|\le B$ and $B=\max(1,\lfloor \rho N_t \rfloor)$.
Selection combines (i)~a \emph{step utility score} $S_k$ derived from external trajectory signals and (ii)~a \emph{token saliency score} $T_i$ from H2O-style attention aggregation.

Concretely, we assign each interaction step a semantic score~$S_k$, fuse it with token-level scores~$T_i$ into a combined score~$C_i$, and evict low-$C_i$ positions.
The pipeline has four stages at every step boundary:
(1)~step span finalization;
(2)~Phase~1 initialization of the newly finalized step score;
(3)~Phase~2 global citation update over all prior steps; and
(4)~Phase~3 joint scoring followed by two-stage budgeted retention over $C_i$.

Figure~\ref{fig:method} gives an overview; details follow.

\begin{figure}[t]
\centering
\IfFileExists{assets/stepkv_method_main.pdf}{%
  \includegraphics[width=\linewidth]{assets/stepkv_method_main.pdf}%
}{%
  \fbox{\parbox[c][0.22\linewidth][c]{0.95\linewidth}{\centering\small Export \texttt{assets/stepkv\_method\_main.drawio} to PDF.}}%
}
\caption{StepKV overview. Each step boundary runs span finalization, three-phase scoring ($S_k$ init, citation update, $C_i$ computation), and budgeted cache retention.}
\label{fig:method}
\end{figure}

\subsection{Step Definition and Finalization}
\label{sec:step-def}

\paragraph{Cache layout.}
The KV cache is split into a protected prefix and a prunable trajectory:
\begin{equation}
\mathcal{C}^{(t)} = \mathcal{C}_{\mathrm{prot}} \;\|\; \mathcal{C}_{\mathrm{traj}}^{(t)},
\end{equation}
where $\mathcal{C}_{\mathrm{prot}}$ stores the system prompt, in-context demonstrations, and user question, and is never pruned.
All Thought/Action/Observation tokens appended during the episode belong to $\mathcal{C}_{\mathrm{traj}}^{(t)}$ and are eligible for eviction.

\paragraph{Step content.}
Step~$k$ comprises the LLM-generated Thought/Action tokens and the environment Observation prefilled afterward:
\begin{equation}
    \text{Step}_k = [\text{Thought}_k, \text{Action}_k \;\|\; \text{Obs}_k],
\end{equation}
mapped to a contiguous global token span $[s_k, e_k]$ in $\mathcal{C}_{\mathrm{traj}}$.
Thought and action tokens are \emph{decoded}; observation tokens are \emph{prefilled} from the environment after the action is executed (e.g., retrieved Wikipedia passages).

\paragraph{Finalization timing.}
We \emph{finalize} step~$k$ at the beginning of step~$k{+}1$, when $\text{Obs}_k$ has been appended and the span $[s_k,e_k]$ is complete.
Until finalization, step~$k$'s observation boundary is unknown, so no step score or pruning span is registered for it.
Only after finalization do we:
\begin{itemize}[leftmargin=*, nosep]
    \item extract anchors $\mathcal{A}_k$ from $\text{Obs}_k$;
    \item compute progression reward $r_k$ and initialize $S_k$ (Phase~1);
    \item register $[s_k,e_k]$ in the step-span table used by Phase~3 retention.
\end{itemize}
This ensures that novelty and success signals are computed on the full observation, not on partial Thought/Action prefixes.

\paragraph{Pruning schedule.}
StepKV runs at step boundaries: after Phase~1--2 scoring and before decoding the next Thought/Action block, Phase~3 may compress $\mathcal{C}_{\mathrm{traj}}^{(t)}$ to budget~$B$.
We set \texttt{observation\_window=0}, so recent observations are not unconditionally protected (unlike H2O baselines with a trailing window of 32 tokens), allowing StepKV to prioritize globally useful observations via $S_k$ rather than recency alone.

\subsection{Anchor Extraction}
\label{sec:anchor}

Anchors are lightweight lexical cues used for novelty and citation.
They require no extra model calls and are extracted deterministically from text.

\paragraph{Extraction procedure.}
From each observation and from the current-step input text used for citation, we build $\mathcal{A}(\text{text})$ as follows:
\begin{enumerate}[leftmargin=*, nosep]
    \item \textbf{Regex scan.} Lowercase the text and match two token types in left-to-right order:
    \begin{itemize}[leftmargin=*, nosep]
        \item \emph{word anchors:} $[a-z][a-z0-9_-]{3,}$ (letter-led alphanumeric tokens);
        \item \emph{numeric anchors:} $\d+$ (every digit sequence, with \textbf{no} length filter).
    \end{itemize}
    \item \textbf{Filtering.} For word anchors only, discard tokens with length $\le 3$ and tokens in a fixed stopword list (e.g., \texttt{thought}, \texttt{action}, \texttt{observation}, \texttt{search}, \texttt{lookup}, \texttt{could}, \texttt{invalid}).
    Numeric anchors bypass the length and stopword filters.
    \item \textbf{Deduplication and cap.} Keep first occurrences in order; retain at most $n=64$ distinct anchors:
    \begin{equation}
        \mathcal{A}(\text{text}) = \text{Top-}n\!\left(\text{dedup}\big(\text{filter}(\text{regex}(\text{text}))\big)\right).
    \end{equation}
\end{enumerate}

\paragraph{Roles in scoring.}
Anchors support two signals:
\begin{itemize}[leftmargin=*, nosep]
    \item \textbf{Novelty} ($\mathrm{nov}_k$): fraction of anchors in $\mathcal{A}_k=\mathcal{A}(\text{Obs}_k)$ that have not appeared in any earlier observation.
    \item \textbf{Citation} ($c_k$): accumulated evidence that later step text re-mentions anchors from step~$k$.
\end{itemize}
Each finalized step stores $\mathcal{A}_k$ in metadata together with $(r_k,c_k,S_k)$.

\subsection{Three-Phase Scoring and Cache Retention}
\label{sec:three-phase}

At each step boundary, StepKV scores the cache in three phases before retention.
For every finalized step $k\in\mathcal{S}^{(t)}$, we maintain anchor set $\mathcal{A}_k$, progression reward $r_k$, citation count $c_k\ge 0$, and step score $S_k$.
All step scores share the same mapping
\begin{equation}
\label{eq:step-score}
f(r,c)=\operatorname{clip}\!\left(w_r\,\operatorname{clip}(r,-1,2)+w_c\log(1+c),\,0,\,S_{\max}\right),
\end{equation}
with defaults $w_r=0.85$, $w_c=0.15$, and $S_{\max}=8$.
Clipping $r$ to $[-1,2]$ prevents a single failed or overly repetitive step from dominating retention; $\log(1+c)$ saturates citation growth so delayed reuse helps without overwhelming token saliency.

\paragraph{Phase 1: New step score initialization.}
When step $k$ is finalized at the start of step $k{+}1$:
\begin{enumerate}[leftmargin=*, nosep]
    \item \textbf{Action normalization.} Parse $\text{Action}_k$ into tool name and argument; normalize whitespace/case to obtain a canonical action string $a_k$ (e.g., \texttt{search[entity name]}).
    \item \textbf{Repetition count.} Let $\mathrm{rep}_k$ be the number of times $a_k$ appeared among actions of steps $<k$. Increment the global action counter after reading $\mathrm{rep}_k$.
    \item \textbf{Anchor/novelty.} Extract $\mathcal{A}_k$ from $\text{Obs}_k$, update the global seen-anchor set, and compute
    \[
    \mathrm{nov}_k = \frac{|\mathcal{A}_k\setminus \mathcal{A}^{\mathrm{seen}}_{<k}|}{\max(1,|\mathcal{A}_k|)}.
    \]
    \item \textbf{Success flag.} Set $\mathrm{succ}_k=0$ if $\text{Obs}_k$ contains failure substrings (\texttt{could not find}, \texttt{invalid action}, \texttt{no results}); otherwise $\mathrm{succ}_k=1$.
    \item \textbf{Reward and score.} Compute
    \begin{equation}
    \label{eq:reward}
    r_k = \mathrm{succ}_k + \mathrm{nov}_k - 0.3\cdot\mathrm{rep}_k,
    \end{equation}
    initialize $c_k\leftarrow 0$, and set $S_k = f(r_k, 0)$ (Eq.~\eqref{eq:step-score}).
\end{enumerate}
Repetition penalizes \emph{only} the newly finalized step in Phase~1; older steps are never retroactively penalized for repeated actions taken later.

\paragraph{Phase 2: Global step score update.}
Before generating step~$t$, we scan the \emph{current input text} $x_t$ (the observation about to be prefilled, plus any adjacent context injected at this boundary).
Let $\mathcal{R}_t=\mathcal{A}(x_t)$.
For \emph{every} finalized step $k<t$:
\begin{enumerate}[leftmargin=*, nosep]
    \item compute overlap count $o_{k,t}=|\mathcal{A}_k\cap\mathcal{R}_t|$;
    \item if $o_{k,t}=0$, leave $(c_k,S_k)$ unchanged;
    \item otherwise update citation and refresh the step score:
    \begin{equation}
    \label{eq:cite}
    c_k \leftarrow c_k + \min\!\left(1,\;\eta\cdot\frac{o_{k,t}}{\max\!\left(1,\min(|\mathcal{A}_k|,|\mathcal{R}_t|)\right)}\right),
    \qquad
    S_k \leftarrow f(r_k, c_k),
    \end{equation}
    with $\eta=1.5$.
\end{enumerate}
The overlap ratio uses the smaller anchor-set size as denominator, so partial reuse of a long observation still yields a meaningful increment; the outer $\min(1,\cdot)$ caps each round's contribution.
Because $r_k$ is fixed after Phase~1, citation exclusively drives upward revisions of older step scores when their entities reappear---implementing delayed-reuse detection without any attention lookahead.

\paragraph{Phase 3: Joint token--step scoring and cache retention.}
Phase~3 assigns a combined score to every prunable token and selects $\mathcal{K}^{(t)}$ under budget~$B$.

\subparagraph{Token saliency ($T_i$).}
For each prunable index $i\in\mathcal{T}^{(t)}$, we compute an H2O heavy-hitter score from piggyback attentions collected during the latest observation prefill.
Let $\mathrm{Attn}^{(\ell)}_{h,q,j}$ be attention from query position $q$ and head $h$ in layer $\ell$ to KV position $j$.
Using the last $L=3$ layers, we define
\begin{equation}
\label{eq:h2o}
T_i = \frac{1}{L}\sum_{\ell\in\mathcal{L}}\sum_{h,q}\mathrm{Attn}^{(\ell)}_{h,q,i},
\end{equation}
where $\mathcal{L}$ indexes the final three decoder layers.
Piggyback scoring reuses the prefill forward pass, avoiding an extra scoring forward pass at each step boundary.

\subparagraph{Combined score ($C_i$).}
Each token inherits the step score of its span label $\mathrm{step}(i)\in\mathcal{S}^{(t)}$:
\begin{equation}
\label{eq:combined}
C_i = \alpha T_i + \beta S_{\mathrm{step}(i)},
\end{equation}
with defaults $\alpha=\beta=0.8$.
Tokens in the same step share the same $S_{\mathrm{step}(i)}$ but may differ in $T_i$, allowing locally salient tokens to survive even within a low-utility step, and high-utility steps to lift their entire span.

\subparagraph{Budget.}
Let $n=N_t=|\mathcal{T}^{(t)}|$ and keep ratio $\rho$.
The global retention budget is
\begin{equation}
\label{eq:budget}
B = \max\!\left(1, \lfloor \rho n \rfloor\right).
\end{equation}

\subparagraph{Two-stage retention within Phase~3.}
Although Phase~1--2 operate at step granularity, eviction is token-level and may leave a step under-represented if all its tokens have low current attention.
We therefore retain tokens in two sub-stages, both ranking by $C_i$:
\begin{enumerate}[leftmargin=*, nosep]
    \item \textbf{Stage 1 (per-step floor).}
    For each step $k\in\mathcal{S}^{(t)}$, let $\mathcal{I}_k=\{i\in\mathcal{T}^{(t)}:\mathrm{step}(i)=k\}$ be its current token indices.
    Keep the top
    \[
    k_k = \max\!\left(1,\;\left\lceil 0.10\cdot|\mathcal{I}_k|\right\rceil\right)
    \]
    tokens in $\mathcal{I}_k$ with highest $C_i$, forming seed set $\mathcal{K}^{(t)}_{\mathrm{floor}}$.
    \item \textbf{Stage 2 (global fill).}
    Starting from $\mathcal{K}^{(t)}_{\mathrm{floor}}$, add tokens from $\mathcal{T}^{(t)}\setminus\mathcal{K}^{(t)}_{\mathrm{floor}}$ in descending $C_i$ order until $|\mathcal{K}^{(t)}|=B$.
    If Stage~1 already selects more than $B$ tokens, trim to the top-$B$ by $C_i$ globally.
\end{enumerate}
The final cache keeps the highest-$C_i$ positions while enforcing both a per-step presence floor and the global length target~$B$.
Note that $S_k$ for old steps does not decrease over time; tokens from those steps may still be evicted because $T_i$ is recomputed every step and longer observations increase $n$, tightening competition under fixed~$\rho$.

\begin{table}[t]
\centering
\caption{StepKV hyperparameters (defaults).}
\label{tab:stepkv-hparams}
\small
\begin{tabular}{@{}lll@{}}
\toprule
Symbol / parameter & Default & Role \\
\midrule
$|\mathcal{A}|$ cap & 64 & max anchors per text segment \\
repeat coeff. & $0.3$ & penalty weight in $r_k$ (Eq.~\eqref{eq:reward}) \\
$w_r, w_c$ & $0.85, 0.15$ & reward / citation weights in Eq.~\eqref{eq:step-score} \\
cite increment cap & $\min(1, 1.5 \times \text{overlap ratio})$ & per-round update in Eq.~\eqref{eq:cite} \\
$\alpha, \beta$ & $0.8, 0.8$ & token / step weights in Eq.~\eqref{eq:combined} \\
$\rho$ & $0.2$--$0.5$ & global KV keep ratio (Eq.~\eqref{eq:budget}) \\
Stage-1 span ratio & $10\%$ & per-step minimum retention in Phase~3 \\
$L$ & $3$ & last layers for H2O score $T_i$ (Eq.~\eqref{eq:h2o}) \\
\bottomrule
\end{tabular}
\end{table}

\begin{algorithm}[t]
\caption{StepKV cache selection at step boundary $t$}
\label{alg:stepkv}
\begin{algorithmic}[1]
\Require Prunable indices $\mathcal{T}^{(t)}$, finalized steps $\mathcal{S}^{(t)}$, incoming text $x_t$, keep ratio $\rho$
\Ensure Selected set $\mathcal{K}^{(t)}$ with $|\mathcal{K}^{(t)}|\le B$
\State \textbf{Phase 1 (init):} finalize step $t{-}1$; compute $r_{t-1}$ via Eq.~\eqref{eq:reward}; set $c_{t-1}\leftarrow 0$; $S_{t-1}\leftarrow f(r_{t-1},0)$
\State \textbf{Phase 2 (cite):} $\mathcal{R}_t\leftarrow\mathcal{A}(x_t)$
\For{$k\in\mathcal{S}^{(t)}$}
    \If{$|\mathcal{A}_k\cap\mathcal{R}_t|>0$}
        \State update $c_k,S_k$ via Eq.~\eqref{eq:cite}
    \EndIf
\EndFor
\State \textbf{Phase 3 (score):} compute $T_i$ for all $i\in\mathcal{T}^{(t)}$ via Eq.~\eqref{eq:h2o}; set $C_i=\alpha T_i+\beta S_{\mathrm{step}(i)}$
\State $B\leftarrow\max(1,\lfloor \rho|\mathcal{T}^{(t)}|\rfloor)$
\State \textbf{Phase 3 (Stage 1):} for each $k\in\mathcal{S}^{(t)}$, add top-$\lceil 0.10|\mathcal{I}_k|\rceil$ tokens by $C_i$ to $\mathcal{K}^{(t)}_{\mathrm{floor}}$
\State \textbf{Phase 3 (Stage 2):} fill $\mathcal{K}^{(t)}$ from $\mathcal{T}^{(t)}\setminus\mathcal{K}^{(t)}_{\mathrm{floor}}$ by descending $C_i$ until $|\mathcal{K}^{(t)}|=B$ (trim if Stage~1 exceeds $B$)
\State \Return $\mathcal{K}^{(t)}$
\end{algorithmic}
\end{algorithm}

\section{Experiments}

\subsection{Experimental Setup}

\paragraph{Models.}
Our main results use \textbf{Qwen2.5-7B-Instruct}; we additionally report Llama-3.1-8B-Instruct results when available.

\paragraph{Datasets and agent protocol.}
We evaluate on multi-step QA benchmarks commonly used with ReAct agents:
\textbf{HotpotQA}, \textbf{2WikiMultihopQA}, \textbf{MuSiQue}, and \textbf{BrowseComp}.
Each dataset uses 500 development samples unless noted otherwise (BrowseComp: 100--518 samples), with fixed random seed $233$.
The agent follows a 6-shot ReAct prompt, uses BM25 retrieval over full Wikipedia ($\sim$5.2M articles, top-5 passages), and is limited to 7 reasoning steps.

\paragraph{Baselines.}
We compare against:
(1)~\textbf{ReAct} without KV pruning;
(2)~\textbf{Full KV} (\texttt{react\_kv\_none});
(3) token-level \textbf{H2O} and \textbf{SnapKV};
(4)~\textbf{TOVA};
(5) step-protected \textbf{step\_anchor\_h2o};
(6) attention-EMA step baseline \textbf{step\_inter};
and (7)~\textbf{StepKV} (\texttt{step\_aware\_h2o}).

\paragraph{Metrics.}
We report Exact Match (EM) and token-level F1.
We also log peak GPU memory (excluding static model weights), average decode cache length after each step, and wall-clock time per sample.

\paragraph{StepKV configuration.}
Unless otherwise noted: $\alpha=\beta=0.8$, $w_r=0.85$, $w_c=0.15$, $\lambda=0.3$, $\eta=1.5$, $L=3$, \texttt{observation\_window=0}, \texttt{attn\_mode=piggyback}.
KV budgets are reported as cache ratios $\rho=50\%$ and $\rho=35\%$ (HotpotQA) or $\rho=20\%$ where indicated.

\subsection{Main Results}

Table~\ref{tab:main-results} summarizes accuracy under matched KV budgets.
On HotpotQA with Qwen2.5-7B at 50\% budget, StepKV achieves \textbf{26.50 EM / 35.50 F1}, closely approaching Full KV (28.40 / 38.11) and ReAct (28.00 / 39.19), while H2O and SnapKV collapse to 12.80/20.24 and 14.30/20.12 respectively.
At a tighter 35\% budget, StepKV retains 24.33 EM and 33.46 F1, whereas H2O falls to 10.50 EM.

On 2WikiMultihopQA at 50\% budget, StepKV reaches 23.80 EM and 28.35 F1 vs.\ 12.20/15.80 for H2O and 33.80/39.10 for Full KV.
The pattern holds across datasets: StepKV consistently recovers most of the full-cache accuracy while token-only methods degrade sharply under the same memory constraint.

\begin{table*}[!t]
\centering
\setlength{\tabcolsep}{4pt}
\resizebox{\linewidth}{!}{%
\begin{tabular}{ll ccc ccc ccc}
\toprule
& & \multicolumn{3}{c}{HotpotQA} & \multicolumn{3}{c}{2WikiMultihopQA} & \multicolumn{3}{c}{BrowseComp} \\
\cmidrule(lr){3-5} \cmidrule(lr){6-8} \cmidrule(lr){9-11}
Model & Method & EM & F1 & KV Budget & EM & F1 & KV Budget & EM & F1 & KV Budget \\
\midrule
\multirow{8}{*}{Qwen2.5-7B}
& Full KV & 28.40 & 38.11 & 100\% & 33.80 & 39.10 & 100\% & -- & -- & -- \\
& ReAct & 28.00 & 39.19 & -- & 32.60 & 38.92 & -- & -- & -- & -- \\
& H2O & 12.80 & 20.24 & 50\% & 12.20 & 15.80 & 50\% & -- & -- & -- \\
& H2O & 10.50 & 17.38 & 35\% & -- & -- & -- & -- & -- & -- \\
& SnapKV & 14.30 & 20.12 & 50\% & -- & -- & -- & -- & -- & -- \\
& SnapKV & 11.38 & 18.67 & 35\% & -- & -- & -- & -- & -- & -- \\
& StepKV (Ours) & \textbf{26.50} & \textbf{35.50} & 50\% & \textbf{23.80} & \textbf{28.35} & 50\% & -- & -- & -- \\
& StepKV (Ours) & \textbf{24.33} & \textbf{33.46} & 35\% & -- & -- & -- & -- & -- & -- \\
\midrule
\multirow{8}{*}{Llama-3.1-8B}
& Full KV & -- & -- & -- & -- & -- & -- & -- & -- & -- \\
& ReAct & -- & -- & -- & -- & -- & -- & -- & -- & -- \\
& H2O & -- & -- & -- & -- & -- & -- & -- & -- & -- \\
& H2O & -- & -- & -- & -- & -- & -- & -- & -- & -- \\
& SnapKV & -- & -- & -- & -- & -- & -- & -- & -- & -- \\
& SnapKV & -- & -- & -- & -- & -- & -- & -- & -- & -- \\
& StepKV (Ours) & -- & -- & -- & -- & -- & -- & -- & -- & -- \\
& StepKV (Ours) & -- & -- & -- & -- & -- & -- & -- & -- & -- \\
\bottomrule
\end{tabular}}
\caption{Main results on multi-step QA benchmarks under different KV cache budgets. ``KV Budget'' is the fraction of prunable trajectory tokens retained after pruning ($\rho$). ReAct denotes the non-pruning agent baseline. Best pruning results are bold.}
\label{tab:main-results}
\end{table*}

\subsection{Efficiency Analysis}

StepKV targets the accuracy--memory trade-off rather than raw speed.
Because hybrid scoring adds lightweight anchor extraction and step metadata updates, wall-clock time is comparable to H2O when using piggyback attention scoring.
However, peak decode cache length and GPU memory track the configured budget $\rho$ similarly to other pruning methods.
Table~\ref{tab:efficiency} reports average peak memory and final cache length; StepKV reduces memory relative to Full KV while retaining far higher EM than H2O at the same budget.

\begin{table}[t]
\centering
\setlength{\tabcolsep}{5pt}
\begin{tabular}{lccc}
\toprule
Method & Peak Mem. (MB) $\downarrow$ & Avg.\ Cache Len $\downarrow$ & Hotpot EM/F1 \\
\midrule
Full KV & -- & -- & 28.40 / 38.11 \\
H2O (50\%) & -- & -- & 12.80 / 20.24 \\
StepKV (50\%) & -- & -- & 26.50 / 35.50 \\
\bottomrule
\end{tabular}
\caption{Efficiency statistics on HotpotQA (Qwen2.5-7B). Fill from experiment logs (\texttt{peak\_memory\_mb}, \texttt{step\_remaining\_tokens}).}
\label{tab:efficiency}
\end{table}

\subsection{Ablation Studies}

\paragraph{Mixing weights ($\beta$ sweep).}
We vary $\beta\in\{1.0,0.8,0.6,0.4,0.2\}$ with $\alpha=1-\beta$ on HotpotQA and 2WikiMultihopQA.
Performance peaks in a mixed regime rather than at pure token-only ($\beta=0$) or pure step-only ($\beta=1$) scoring, confirming that $T_i$ and $U_s$ are complementary.

\paragraph{External step utility vs.\ attention EMA (\texttt{step\_inter}).}
Replacing external reward/citation utility with an attention exponential moving average over step spans reduces accuracy and interpretability.
External signals are cheaper and directly encode delayed reuse through citation.

\paragraph{Citation and progression signals.}
Removing citation updates (Phase~2) or novelty terms in $R_s$ degrades performance, especially on longer trajectories where reuse occurs several steps after an observation.
Forced removal of the highest-utility step (\texttt{force\_drop\_step\_ids}) hurts EM more than removing a random or lowest-utility step, validating that $U_s$ tracks semantically important reasoning stages.

\subsection{Qualitative Analysis}

Figure~\ref{fig:case} illustrates delayed reuse.
In a successful StepKV trajectory, an early observation introduces rare entities that receive moderate attention initially.
Several steps later, the model cites the same anchors in a new search or thought; Phase~2 raises the earlier step utility and Phase~3 retains its tokens despite low current $T_i$.
Under H2O, the same tokens are evicted earlier and the agent fails to recover the correct answer.

\begin{figure}[t]
\centering
\IfFileExists{assets/case_stepkv_step_intervals.pdf}{%
  \includegraphics[width=\linewidth]{assets/case_stepkv_step_intervals.pdf}%
}{%
  \fbox{\parbox[c][0.22\linewidth][c]{0.95\linewidth}{\centering\small Export \texttt{assets/case\_stepkv\_step\_intervals.drawio} to PDF.}}%
}
\caption{Case study of delayed reuse. Step intervals are shown along the trajectory; citation arrows indicate Phase~2 updates that increase utility of earlier steps when later text reuses their anchors.}
\label{fig:case}
\end{figure}

\paragraph{Why old steps still lose tokens without $U_s$ decreasing.}
Step utility is monotone non-decreasing under our update rules: repetition does not retroactively lower older steps, and citation only increases $C_s$.
Nevertheless, eviction can remove old step tokens because $T_i$ is recomputed every step and long new observations increase $n$, so a fixed ratio $\rho$ forces global competition among all surviving tokens via $C_i=\alpha T_i+\beta U_s$.

\section{Conclusion}

We revisit KV cache compression for long-horizon LLM agents and show that token-level methods fail because importance is step-conditional and often delayed across ReAct trajectories.
StepKV addresses this by grouping tokens into reasoning steps, estimating step utility from lightweight external signals, and performing three-phase scoring at each step boundary before hybrid retention with H2O token scores.
Across multi-hop QA benchmarks, StepKV substantially outperforms H2O and SnapKV under matched budgets and approaches full-cache accuracy.
We hope this step-aware perspective encourages future KV systems to exploit trajectory structure rather than relying solely on local attention heuristics.

\section*{Limitations}

Our step utility relies on simple anchor heuristics rather than learned reward models; noisy tool outputs or non-English text may reduce signal quality.
Results are primarily reported at 7B--8B scale; larger models and additional agent environments (e.g., web browsing with HTML observations) require further validation.
Anchor rules do not capture structured tables or code snippets comprehensively.
Finally, although the paper presents a three-phase scoring narrative, the reference implementation adds an optional per-step retention floor before global top-$B$ selection (Appendix~\ref{app:implementation}); future work may unify these into a single constrained optimization rule.

\section*{Acknowledgments}

Anonymous for review.

\appendix

\section{Implementation Details of Budget Retention}
\label{app:implementation}

After computing hybrid scores $C_i$, our reference implementation performs optional \emph{pool-wise} retention before global top-$B$ fill:
for each finalized step span, at least $\lceil 0.10 \cdot |\mathcal{I}_s|\rceil$ tokens with the highest $C_i$ within that span are reserved, then the remaining budget is filled globally by descending $C_i$.
This prevents any single step from being entirely wiped out before global competition.
When \texttt{step\_poolwise\_prune=False}, a floor based on \texttt{step\_aware\_min\_keep} and \texttt{step\_aware\_min\_keep\_ratio} is used instead.

\section{Hyperparameters}
\label{app:hyperparams}

\begin{table}[h]
\centering
\small
\begin{tabular}{ll}
\toprule
Parameter & Value \\
\midrule
Cache ratio $\rho$ & 0.50, 0.35, 0.20 \\
$\alpha$, $\beta$ & 0.8, 0.8 \\
$w_r$, $w_c$ & 0.85, 0.15 \\
$\lambda$ (repeat penalty) & 0.3 \\
$\eta$ (citation scale) & 1.5 \\
$U_{\max}$ & 8 \\
Attention scoring layers $L$ & 3 \\
Max ReAct steps & 7 \\
BM25 top-$k$ & 5 \\
Random seed & 233 \\
Observation window & 0 (StepKV) / 32 (H2O baseline) \\
Attention mode & piggyback (StepKV) \\
\bottomrule
\end{tabular}
\caption{Default experimental hyperparameters.}
\label{tab:hyperparams}
\end{table}

\clearpage
\bibliography{custom}

\end{document}
